\input{../tex-inputs/latex-leseansicht-vorspann}

\section[Arthur, Olga Schnitzler und Elisabeth Steinrück an Gustav Schwarzkopf, 14. 11. 1911]{L04521 Arthur, Olga Schnitzler und Elisabeth Steinrück an Gustav Schwarzkopf, 14. 11. 1911}
\nopagebreak\mylabel{L04521v}
\rehead{ }\normalsize\beginnumbering\briefempfaengerindex{Schwarzkopf, Gustav@\textsc{Schwarzkopf, Gustav}!zzzSteinrück, Elisabeth@\emph{von Elisabeth Steinrück}!1911-11-142@{14. 11. 1911}|(be}\briefempfaengerindex{Schwarzkopf, Gustav@\textsc{Schwarzkopf, Gustav}!zzzSchnitzler, Olga@\emph{von Olga Schnitzler}!1911-11-142@{14. 11. 1911}|(be}\briefempfaengerindex{Schwarzkopf, Gustav@\textsc{Schwarzkopf, Gustav}!zzzSchnitzler, Arthur@\emph{von Arthur Schnitzler}!1911-11-142@{14. 11. 1911}|(be}
\toendnotes[C]{\smallbreak\pagebreak[2]}
\correspDesc{Versand  durch Arthur Schnitzler, Olga Schnitzler, Elisabeth Steinrück am 14. 11. 1911 in Partenkirchen
\newline{}Übermittlung  am 14. 11. 1911 in Partenkirchen
\newline{}Erhalt  durch Gustav Schwarzkopf im Zeitraum [15. 11. 1911 – 19. 11. 1911?] in Wien}\toendnotes[C]{\smallbreak}
\Standort{DLA, A:Schnitzler, HS.1985.1.1897.}
\physDesc{Bildpostkarte, 212 Zeichen
\newline{}Handschrift Arthur Schnitzler: Bleistift, deutsche Kurrent
\newline{}Handschrift Olga Schnitzler: Bleistift, lateinische Kurrent
\newline{}Handschrift Elisabeth Steinrück: Bleistift, lateinische Kurrent
\newline{}Versand: Stempel: »\nobreak{}\oindex{Partenkirchen@\textbf{Partenkirchen}|pwk}Partenkirchen 1, 14 1\textcolor{gray}{1}{ }{[}11{]}, \textcolor{gray}{VII}\nobreak{}«.  }
\pstart
           \noindent{}\centering{}{\pb}\textcolor{gray}{\textbf{Partenkirchen\oindex{Garmisch-Partenkirchen@\textbf{Garmisch-Partenkirchen}|pw}
                  mit Zugspitze\oindex{Zugspitze@\textbf{Zugspitze}|pw}.}}\pend
           \pstart{}{\pb}Herrn\pend{}\pstart{}Guſtav Schwarzkopf\pend{}\pstart{}Wien I\oindex{I., Innere Stadt@\textbf{I., Innere Stadt}|pw}\pend{}\pstart{}\textsc{Tiefer Graben} 17\oindex{Wien@\textbf{Wien}!I., Innere Stadt@\textbf{I., Innere Stadt}!Tiefer Graben 17@\textbf{Tiefer Graben 17}|pw}\pend{}{\bigskip}\vspace{1em}
\pstart
           \noindent{}Herzlichſt Ihr\pend
           \pstart \spacefill\mbox{Arthur}\pend{}\selectlanguage{ngerman}\vspace{1em}
\pstart
           \noindent{}{[}hs. O. Schnitzler:{]} Das Lisl\pwindex{Steinrück, Elisabeth 19.\,11.\,1885 – 7.\,4.\,1920 Partenkirchen@\textsc{Steinrück, Elisabeth} (19.\,11.\,1885 – 7.\,4.\,1920 Partenkirchen)|pw} ist dicker
               geworden und sieht so gut aus, wie
               seit langem nicht. Gott sei Dank.\pend
           
\pstart
           Herzliche Grüsse{\\[\baselineskip]}\spacefill\mbox{Olga Schnitzler}\pend
           \leftskip=0em{}\selectlanguage{ngerman}\vspace{1em}
\pstart
           \noindent{}{[}hs. Steinrück:{]} Und Ihre alte\pend
           \pstart \spacefill\mbox{Lise St.}\pend{}\selectlanguage{ngerman}\vspace{1em}
\pstart
           {[}hs. A. Schnitzler:{]} {\pb}14.\,11.\,911\pend
           \selectlanguage{ngerman}\endnumbering\briefempfaengerindex{Schwarzkopf, Gustav@\textsc{Schwarzkopf, Gustav}!zzzSteinrück, Elisabeth@\emph{von Elisabeth Steinrück}!1911-11-142@{14. 11. 1911}|)be}\briefempfaengerindex{Schwarzkopf, Gustav@\textsc{Schwarzkopf, Gustav}!zzzSchnitzler, Olga@\emph{von Olga Schnitzler}!1911-11-142@{14. 11. 1911}|)be}\briefempfaengerindex{Schwarzkopf, Gustav@\textsc{Schwarzkopf, Gustav}!zzzSchnitzler, Arthur@\emph{von Arthur Schnitzler}!1911-11-142@{14. 11. 1911}|)be}\mylabel{L04521h}
\begin{anhang}
\end{anhang}\newcommand{\dateiname}{L04521}\newcommand{\titel}{Arthur, Olga Schnitzler und Elisabeth Steinrück an Gustav Schwarzkopf, 14. 11. 1911}\newcommand{\editorInnen}{Herausgegeben von Jahnke, SelmaMüller, Martin Anton}\input{../tex-inputs/latex-leseansicht-abspann}
